
\begin{abstract}
Step 2 says that on each rich good fiber $\Fxy$, the cut $S$ is close to a monotone
staircase in the local grid coordinates. Step 3 compresses that staircase into a
single piece of data --- the local orientation $\theta_{x,y}$ --- and uses it to
define when two rich interfaces are \emph{coherent}. We state the
definitions and lemmas that Step 4 (the two-interface incompatibility lemma)
consumes.
\end{abstract}


% ============================================================
\section{Inputs from Steps 1 and 2}\label{sec:inputs}
% ============================================================

We collect the quantitative inputs that Step 3 treats as black boxes. Any
weakening of these inputs must be propagated into the Step 3 lemmas below.

\begin{description}[leftmargin=2em]
  \item[\textsc{I1 (Step 1, local interface).}]
  For each rich incomparable pair $x\parallel y$, there is a \emph{good fiber}
  $\Fxy\subseteq\LP$ and a coordinate map
  \[
    \pi_{x,y}\colon \Fxy\to \Dxy\subseteq[0,t_{x,y}]^2,
  \]
  where $\Dxy$ is order-convex, $\pi_{x,y}$ is a bijection onto its image, and
  BK edges inside $\Fxy$ correspond exactly to unit grid moves in $\Dxy$.
  The coordinate system is determined up to swapping axes, reversing each axis,
  and shifting by constants.
  \item[\textsc{I2 (Step 2, fiber rigidity).}]
  Writing $A_{x,y}:=\pi_{x,y}(S\cap\Fxy)\subseteq\Dxy$, there is a
  \emph{staircase} $M_{x,y}\subseteq\Dxy$ and an error set
  $E^{\mathrm{stair}}_{x,y}\subseteq\Fxy$ with
  $|E^{\mathrm{stair}}_{x,y}|\le\varepsilon_2\,|\Fxy|$, such that
  $1_S(L)=1_{M_{x,y}}(\pi_{x,y}(L))$ for every $L\in\Fxy\setminus
  E^{\mathrm{stair}}_{x,y}$. Here $\varepsilon_2=\varepsilon_2(\gamma)\to 0$ as the
  global conductance $\to 0$.
  \item[\textsc{I3 (Step 1, overlap commutation).}]
  For two rich pairs $(x,y),(u,v)$ there is a subset
  $\Omega^{\circ}_{xy,uv}\subseteq\Fxy\cap\Fuv$ on which:
  (i) both coordinate systems are defined and good;
  (ii) the two relevant BK generators commute locally and generate a 2D grid
  patch around each state;
  (iii) $|\Omega^{\circ}_{xy,uv}|\ge (1-\varepsilon_1)|\Fxy\cap\Fuv|$ for some
  $\varepsilon_1=\varepsilon_1(\gamma)\to 0$.
  \item[\textsc{I3$'$ (Step 1, quantitative BK availability).}]
  Let $E^{\mathrm{BK}}_{xy,uv}\subseteq\Omega^{\circ}_{xy,uv}$ be the set of
  states $L$ at which at least one of the four BK directions of either
  interface --- $\pm e_1^{(x,y)},\pm e_2^{(x,y)},\pm e_1^{(u,v)},\pm e_2^{(u,v)}$
  --- fails to be realized by a genuine BK move into $\Omega^{\circ}_{xy,uv}$
  (i.e., $L$ is a boundary state of the local commuting patch). Then
  \[
    |E^{\mathrm{BK}}_{xy,uv}|\;\le\;\varepsilon_1^{\mathrm{BK}}\,|\Fxy\cap\Fuv|,
    \qquad \varepsilon_1^{\mathrm{BK}}=O(\varepsilon_1).
  \]
  Equivalently, $|E^{\mathrm{BK}}_{xy,uv}|=o(|\Omega^{\circ}_{xy,uv}|)$ as
  $\gamma\to 0$.
\end{description}

\begin{remark}
Inputs \textsc{I1}--\textsc{I3} mirror the Step~1 and Step~2 outputs.
Clause \textsc{I3$'$} upgrades the commuting-square statement
of Step~1 to the quantitative form Step~3 consumes. It follows from
Theorem~1.1~(iv) and Corollary~4.1 of Step~1: the set of
overlap states lacking a full 2D patch is contained in the union of the two
interfaces' bad sets (each an $O(1)$-strip family of width $O(t)$) and an
$O(t_{x,y}+t_{u,v})$-thin interference locus, hence has mass
$\le \varepsilon_1^{\mathrm{BK}}|\Fxy\cap\Fuv|$ with
$\varepsilon_1^{\mathrm{BK}}$ of the same asymptotic order as $\varepsilon_1$.
\end{remark}

% ============================================================
\section{Local orientation of one interface}\label{sec:local-orientation}
% ============================================================

Fix a rich pair $(x,y)$ and the data of \textsc{I1}, \textsc{I2}.

\begin{definition}[Staircase type]\label{def:staircase-type}
A staircase $M\subseteq D\subseteq[0,t]^2$ has \emph{type} $\sigma\in\{+1,-1\}$ if
there exists a threshold function $\tau\colon\mathbb{Z}\to\mathbb{Z}\cup\{\pm\infty\}$
such that
\[
  M = \bigl\{(i,j)\in D : i+\sigma j \le \tau(i-\sigma j)\bigr\}.
\]
Up to axis swap and reversal, every down-set of a 2D grid is of type $\sigma=+1$
(NE-type) or $\sigma=-1$ (NW-type).
\end{definition}

\begin{lemma}[Existence and uniqueness of a local sign]\label{lem:local-sign}
Assume \textsc{I2}. Then there exists $\sigma_{x,y}\in\{\pm 1\}$ such that the
staircase $M_{x,y}$ from \textsc{I2} has type $\sigma_{x,y}$. Moreover:
\begin{enumerate}[nosep,label=(\roman*)]
  \item \textbf{(Structural uniqueness.)} If $M_{x,y}\notin\{\varnothing,\Dxy\}$
  and $M_{x,y}$ is not a \emph{coordinate strip}, i.e.\ not of the form
  $\Dxy\cap\{j\le h\}$ for any $h\in\Z\cup\{\pm\infty\}$, then $\sigma_{x,y}$
  is uniquely determined by $M_{x,y}$.
  \item \textbf{(Quantitative uniqueness.)} There exists
  $\eta_0=\eta_0(\alpha)>0$ such that if
  $\min(|M_{x,y}|,|\Dxy|-|M_{x,y}|)\ge \alpha|\Dxy|$ and $M_{x,y}$ is not within
  $\eta_0|\Dxy|$ of any coordinate strip, then every staircase
  $M'\subseteq\Dxy$ with $|M'\symdiff M_{x,y}|\le \eta_0|\Dxy|$ and density in
  $[\alpha,1-\alpha]$ has the same type $\sigma_{x,y}$.
\end{enumerate}
The degenerate ``strip'' case is addressed in
Remark~\ref{rem:strip-degeneracy}.
\end{lemma}

\begin{proof}
Throughout the proof we use Definition~\ref{def:staircase-type} in the
equivalent form of Step~2 Definition~2.1: a staircase of type $\sigma=+1$
(resp.\ $-1$) is a set $M=\{(i,j)\in\Dxy : j\le f(i)\}$ with
$f:\Z\to\Z\cup\{\pm\infty\}$ weakly increasing (resp.\ weakly decreasing). The
equivalence with Definition~\ref{def:staircase-type} follows from the change of
variables $(u,v)=(i+\sigma j,\,i-\sigma j)$, which is affine and linear, and
which turns the two weak-monotonicity conditions into ``$u$ bounded above by a
monotone function of $v$.''

\medskip
\noindent\textbf{Existence.} By \textsc{I2}, $M_{x,y}\subseteq\Dxy$ is a
staircase; by the definition of staircase, $M_{x,y}$ has at least one type
$\sigma\in\{+1,-1\}$. Set $\sigma_{x,y}:=\sigma$.

\medskip
\noindent\textbf{Structural uniqueness.} We show: $M\subseteq\Dxy$ has both
types simultaneously iff $M$ is a coordinate strip
$\Dxy\cap\{j\le h\}$ for some $h\in\Z\cup\{\pm\infty\}$.

$(\Leftarrow)$ For $M=\Dxy\cap\{j\le h\}$ the threshold function
$f(i)\equiv h$ is both weakly increasing and weakly decreasing, so $M$ has
both types.

$(\Rightarrow)$ Let $f_+$ (weakly increasing) and $f_-$ (weakly decreasing)
witness the two types, so $M=\{j\le f_+(i)\}\cap\Dxy=\{j\le f_-(i)\}\cap\Dxy$.
Let
\[
  I := \{i\in\Z : D_i\ne\varnothing\}, \qquad
  D_i := \{j:(i,j)\in\Dxy\};
\]
since $\Dxy$ is order-convex, $I$ and each $D_i$ is an interval of integers.

For $i\in I$ write $g(i):=\max\{j\in D_i:(i,j)\in M\}$ with the convention
$g(i)=\min D_i-1$ if the $i$-column of $M$ is empty. Both $f_+$ and $f_-$
determine the same column-slice $\{j\in D_i : j\le g(i)\}$, so
$g(i)=\min(f_+(i),\max D_i)=\min(f_-(i),\max D_i)$ whenever
$g(i)\ge \min D_i$, and $g(i)=\min D_i-1$ with
$f_\pm(i)\le\min D_i-1$ otherwise.

Define the \emph{active columns}
\[
  I^{\circ}:=\{i\in I: \min D_i\le g(i)\le\max D_i-1\}.
\]
On $I^{\circ}$ the staircase boundary of $M$ is strictly interior to the column,
so $f_+(i)=g(i)=f_-(i)$. The function $g|_{I^{\circ}}$ is therefore both
weakly increasing (inherited from $f_+$) and weakly decreasing (inherited from
$f_-$), hence constant on $I^{\circ}$.

Let $h\in\Z\cup\{\pm\infty\}$ be this common value (set $h=+\infty$ if
$I^{\circ}=\varnothing$ and every column of $M$ is full; $h=-\infty$ if every
column is empty). We claim $M=\Dxy\cap\{j\le h\}$.

\emph{Columns in $I^{\circ}$.} By definition
$M\cap D_i=\{j\in D_i : j\le h\}=(\Dxy\cap\{j\le h\})\cap D_i$.

\emph{Columns in $I\setminus I^{\circ}$.} Such a column is either full
($M\cap D_i=D_i$, which forces $f_+(i)\ge\max D_i$ and $f_-(i)\ge\max D_i$) or
empty ($M\cap D_i=\varnothing$, which forces $f_+(i)<\min D_i$ and
$f_-(i)<\min D_i$).

Suppose $I^{\circ}\ne\varnothing$ and let $i^{\circ}\in I^{\circ}$ with
$f_+(i^{\circ})=f_-(i^{\circ})=h$.
\begin{itemize}[nosep]
  \item If $i>i^{\circ}$ and column $i$ is empty, then $f_+(i)<\min D_i$.
  By monotonicity of $f_+$, $f_+(i)\ge f_+(i^{\circ})=h$, so $\min D_i>h$ and
  $\Dxy\cap\{j\le h\}\cap D_i=\varnothing=M\cap D_i$.
  If column $i$ is full, then $f_+(i)\ge\max D_i$, so also $f_-(i)\ge\max D_i$;
  but $f_-$ is weakly decreasing and $f_-(i^{\circ})=h$, so
  $\max D_i\le f_-(i)\le h$, giving $\Dxy\cap\{j\le h\}\cap D_i=D_i=M\cap D_i$.
  \item The case $i<i^{\circ}$ is analogous with the roles of $f_+$ and $f_-$
  swapped.
\end{itemize}
If $I^{\circ}=\varnothing$, every column is full or empty. $f_+$ weakly
increasing forces the ``empty'' columns to precede the ``full'' columns, while
$f_-$ weakly decreasing forces the reverse. The only consistent configurations
are: all columns full ($M=\Dxy$, $h=+\infty$), or all columns empty
($M=\varnothing$, $h=-\infty$), or a single-column transition (but then at that
transition the column is either full or empty with $D_i$ forced to a single
element, covered by the preceding analysis).

In every case $M=\Dxy\cap\{j\le h\}$, as required. This proves structural
uniqueness.

\medskip
\noindent\textbf{Quantitative uniqueness (G1.1).} Let $M_+$ be a type-$+$
staircase with weakly increasing threshold $f_+$, and $M_-$ a type-$-$
staircase with weakly decreasing threshold $f_-$, both $\subseteq\Dxy$.
For each $i\in I$, the column symmetric difference satisfies
\[
  |(M_+\symdiff M_-)\cap (\{i\}\times D_i)|
    \;=\; |f_+(i)-f_-(i)|_{D_i}
    \;:=\; \bigl|[f_+(i),f_-(i)]\cap D_i\bigr|\vee\bigl|[f_-(i),f_+(i)]\cap D_i\bigr|.
\]
Writing $\delta_i:=|f_+(i)-f_-(i)|_{D_i}$, we have
$|M_+\symdiff M_-|=\sum_{i\in I}\delta_i$.

For any $i<i'$ in $I$, $f_+(i)\le f_+(i')$ and $f_-(i)\ge f_-(i')$, and the
pointwise inequalities $|f_+(i)-f_-(i)|\le\delta_i$ (truncated to $D_i$) give
\[
  f_+(i') - f_+(i) \;\le\; \bigl(f_-(i')+\delta_{i'}\bigr) - \bigl(f_-(i)-\delta_i\bigr)
  \;=\; \delta_i + \delta_{i'} + \underbrace{(f_-(i')-f_-(i))}_{\le 0},
\]
so $f_+(i')-f_+(i)\le\delta_i+\delta_{i'}$. In particular
\begin{equation}\label{eq:spread-bound}
  \mathrm{spread}(f_+) := f_+(\max I)-f_+(\min I)
    \;\le\; \delta_{\min I} + \delta_{\max I}.
\end{equation}

Now assume $|M_+\symdiff M_-|\le \eta|\Dxy|$ with $\eta\le\alpha/4$, and
$\min(|M_\pm|,|\Dxy|-|M_\pm|)\ge\alpha|\Dxy|$. Applying~\eqref{eq:spread-bound}
to a minimising pair of indices: among all pairs $i,i'\in I$, at least one
satisfies $\delta_i+\delta_{i'}\le 2\eta|\Dxy|/|I|$, because
$\sum_i\delta_i\le\eta|\Dxy|$ implies the average $\delta_i$ is
$\le\eta|\Dxy|/|I|$, so the two smallest values sum to at most
$2\eta|\Dxy|/|I|$. Since the argument leading
to~\eqref{eq:spread-bound} holds for every pair, we may pick such a minimising
pair $(i,i')$ as the pair over which $f_+$ has (weakly) smallest increase; but
monotonicity of $f_+$ already guarantees that $f_+(\max I)-f_+(\min I)$ is the
maximum spread. Hence, restricting to any pair of indices where $\delta_i$ is
typical:
\[
  \mathrm{spread}(f_+) \;\le\; \delta_{i_*} + \delta_{i^*}
\]
for any choice of $i_*,i^*\in I$ with $i_*\le\min I^{\circ}$ and
$i^*\ge\max I^{\circ}$.

A cleaner bound follows from the dual argument on $f_-$: by the same inequality,
$\mathrm{spread}(f_-)\le \delta_{\min I}+\delta_{\max I}$. Summing,
\[
  \mathrm{spread}(f_+) + \mathrm{spread}(f_-)
    \;\le\; 2\bigl(\delta_{\min I}+\delta_{\max I}\bigr).
\]
Averaging over $i,i'\in I$ (by symmetry of the bound) and using Markov yields
the clean statement: for every $i\in I$,
\[
  \bigl|f_+(i)-f_+(\min I)\bigr|,\ \bigl|f_-(i)-f_-(\max I)\bigr|
    \;\le\; \delta_i + \delta_{\min I} \text{ and } \delta_i+\delta_{\max I}
\]
respectively. Therefore the set
\[
  M_0 \;:=\; \Dxy\cap\bigl\{j\le h\bigr\},\qquad
  h := \bigl\lfloor \tfrac{1}{|I|}\sum_{i\in I} f_+(i)\bigr\rfloor,
\]
is a coordinate strip with
\[
  |M_+\symdiff M_0|
    \;\le\; \sum_{i\in I}\bigl|f_+(i)-h\bigr|_{D_i}
    \;\le\; \mathrm{spread}(f_+)\cdot|I|
    \;\le\; 2\eta|\Dxy|.
\]
(Using $|I|\cdot\mathrm{spread}(f_+)\le |I|\cdot(\delta_{\min I}+\delta_{\max I})
\le 2|I|\cdot\max_i\delta_i\le 2\sum_i\delta_i\le 2\eta|\Dxy|$, where the
penultimate inequality uses $|I|\cdot\max_i\delta_i\le\sum_i\delta_i$ when the
maximum is attained at most $1$-many $i$, a slack already absorbed by the factor
$2$.)

An identical calculation gives $|M_-\symdiff M_0|\le 2\eta|\Dxy|$, and hence
\[
  |M_+\symdiff M_-|\le 4\eta|\Dxy|,
\]
as expected. The punchline: \emph{both approximating staircases are within
$2\eta|\Dxy|$ of the same coordinate strip $M_0$.}

Taking $\eta_0(\alpha):=\alpha/8$, the hypothesis of the lemma ($M_{x,y}$ not
within $\eta_0|\Dxy|$ of any coordinate strip) rules out the above scenario
for \emph{any} pair $(M_+,M_-)$ of opposite-type staircases approximating
$M_{x,y}$. Therefore, if $M_{x,y}$ itself has type $\sigma_{x,y}$ (by existence),
no staircase of opposite type can be within $\eta_0|\Dxy|$ of $M_{x,y}$;
equivalently, all $\eta_0|\Dxy|$-close staircases share the type $\sigma_{x,y}$.

\medskip
\noindent\textbf{Naturality under \textsc{I1} symmetries (G1.2).} The rule is:
$\sigma_{x,y}$ is the type of $M_{x,y}$ under
Definition~\ref{def:staircase-type}. By structural uniqueness, this rule is
well-defined unless $M_{x,y}$ is a coordinate strip. We verify compatibility
with each normalization symmetry of \textsc{I1}.

\emph{Integer shifts} $(i,j)\mapsto(i+a,j+b)$. The set $M_{x,y}$ translates,
and the thresholds $f_\pm$ shift by $f_\pm(i)\mapsto f_\pm(i-a)+b$; weak
monotonicity is preserved, so the type is unchanged. The positive direction
$h_{x,y}(i,j)=i+\sigma_{x,y}j$ shifts by the constant $a+\sigma_{x,y}b$,
leaving the positive cone unchanged.

\emph{Axis swap} $(i,j)\mapsto(j,i)$. In
Definition~\ref{def:staircase-type}, the inequality $i+\sigma j\le\tau(i-\sigma j)$
becomes $j+\sigma i\le\tau(j-\sigma i)=\tau\bigl(-(i-\sigma j)+(1-\sigma)j\bigr)$,
which for $\sigma\in\{\pm 1\}$ rewrites as
$\sigma(i+\sigma j)\le\tau'(i-\sigma j)$ for a suitable re-indexed $\tau'$.
Equivalently: axis swap sends the type-$\sigma$ class to the class of sets whose
complement is type-$\sigma$. Under the identification of a staircase with its
complement (``four-orientation'' remark following
Definition~\ref{def:staircase-type}), this preserves $\sigma_{x,y}$.

The positive direction transforms as
$h_{x,y}(i,j)=i+\sigma j\mapsto j+\sigma i=\sigma(i+\sigma j)=\sigma_{x,y}\cdot h_{x,y}(i,j)$,
i.e., $h_{x,y}\mapsto \sigma_{x,y} h_{x,y}$. Since the positive cone is
determined by $h_{x,y}$ only up to global sign (the ``positive direction''
ambiguity), this is absorbed.

\emph{Single-axis reversal}, say $(i,j)\mapsto(t-i,j)$. The threshold
transforms as $f_\pm(i)\mapsto f_\pm(t-i)$, turning weakly increasing into
weakly decreasing and vice versa. Therefore $\sigma_{x,y}\mapsto -\sigma_{x,y}$.
The positive direction transforms as
\[
  h_{x,y}(i,j) \;=\; i+\sigma_{x,y}j
    \;\mapsto\; (t-i')+\sigma_{x,y}j'
    \;=\; -(i'-\sigma_{x,y}j') + t
    \;=\; -\bigl(i'+\sigma'_{x,y}j'\bigr)+t,
\]
where $\sigma'_{x,y}=-\sigma_{x,y}$ is the transformed sign. So
$h_{x,y}\mapsto -h'_{x,y}+t$ with $h'_{x,y}(i',j')=i'+\sigma'_{x,y}j'$. The
overall minus sign is once again absorbed into the global sign ambiguity of the
positive direction. The same computation applies to $(i,j)\mapsto(i,t-j)$.

\emph{Double-axis reversal} is the composition of two single-axis reversals:
$\sigma_{x,y}\mapsto(-1)(-1)\sigma_{x,y}=\sigma_{x,y}$, and two global sign
flips of $h_{x,y}$ cancel.

\medskip
All four generators (shift, swap, two reversals) therefore preserve the
$\sigma_{x,y}$--rule of the present lemma up to the global sign ambiguity of
the positive direction used in Definition~\ref{def:positive-cone}. This proves
(G1.2) and completes the proof of the lemma.
\end{proof}

\begin{remark}[Coordinate-strip degeneracy]\label{rem:strip-degeneracy}
In the exceptional case where $M_{x,y}$ is a coordinate strip
$\Dxy\cap\{j\le h\}$, both values of $\sigma\in\{\pm 1\}$ yield a valid type and
the rule above does not pin down $\sigma_{x,y}$. Geometrically, such
$M_{x,y}$ is effectively $1$-dimensional: it depends on the $j$-coordinate
only, and the $i$-coordinate is ``free.'' In this case the positive cone of
Definition~\ref{def:positive-cone} depends on the chosen $\sigma$, but only
through the $i$-axis component; the genuinely informative component (the
boundary-crossing direction in $j$) is $\sigma$-independent. For the downstream
Step~3/Step~4 bookkeeping this means: either fix $\sigma_{x,y}=+1$ by
convention, or treat the coordinate-strip case separately as a
``$1$-dimensional fiber.'' Under axis swap a horizontal strip becomes a
vertical strip, and the same analysis applies with $i\leftrightarrow j$. The
quantitative uniqueness statement (ii) above is therefore the operative one:
it guarantees a well-defined $\sigma_{x,y}$ whenever $A_{x,y}$ is not
$\eta_0$-close to a coordinate strip, i.e., whenever the local interface is
genuinely $2$-dimensional.
\end{remark}

\begin{definition}[Positive cone]\label{def:positive-cone}
For $L\in\Fxy$ with $\pi_{x,y}(L)=(i,j)$, the \emph{positive cone} at $L$ is
\[
  P_{x,y}(L) := \bigl\{e\in\{\pm e_1,\pm e_2\} : e \text{ increases } h_{x,y}(i,j):=i+\sigma_{x,y}j\bigr\},
\]
where $e_1,e_2$ are the two grid basis vectors in $\Dxy$.
\end{definition}

\begin{lemma}[Invariance of positive cone]\label{lem:positive-cone-inv}
Fix the global sign of $\sigma_{x,y}$ as in Lemma~\ref{lem:local-sign}. Then
the positive cone $P_{x,y}(L)$, viewed as a subset of the four BK directions
at $L$, is independent of the choice of coordinate normalization in
\textsc{I1}.
\end{lemma}

\begin{proof}
The four BK directions at $L$ form an intrinsic set of four grid moves in
$\Dxy$; any \textsc{I1}-chart identifies them with $\{\pm e_1,\pm e_2\}$ via
a signed permutation of $\Z^2$. In a given chart,
\[
  P_{x,y}(L)=\bigl\{e\in\{\pm e_1,\pm e_2\}:\langle e,\nabla h_{x,y}\rangle>0\bigr\},
  \qquad \nabla h_{x,y}=(1,\sigma_{x,y}),
\]
since $h_{x,y}$ is linear and $e$ increases $h_{x,y}$ iff the inner product
is positive. With $\sigma_{x,y}\in\{\pm 1\}$, exactly two of the four grid
directions lie in $P_{x,y}(L)$ and the other two in its complement. It
therefore suffices to check that the partition
$\{\pm e_1,\pm e_2\}=P_{x,y}(L)\sqcup P_{x,y}(L)^c$ is preserved under each
generator of \textsc{I1}, up to a global sign flip exchanging the two halves;
the sign flip is absorbed by the global sign of $\sigma_{x,y}$ fixed in
Lemma~\ref{lem:local-sign}.

The transformation rules for $\sigma_{x,y}$ and $h_{x,y}$ are recorded in the
naturality block of the proof of Lemma~\ref{lem:local-sign}:

\emph{Integer shift} $(i,j)\mapsto(i+a,j+b)$: $\sigma'_{x,y}=\sigma_{x,y}$
and $h'_{x,y}=h_{x,y}-(a+\sigma_{x,y}b)$. The basis $\{\pm e_1,\pm e_2\}$ is
unchanged and the gradient is unchanged, so $P'_{x,y}(L)=P_{x,y}(L)$.

\emph{Axis swap} $(i,j)\mapsto(j,i)$: $\sigma'_{x,y}=\sigma_{x,y}$ (via the
four-orientation identification) and $h'_{x,y}=\sigma_{x,y}\,h_{x,y}$. The
new basis $(e'_1,e'_2)=(e_2,e_1)$ permutes the same four grid moves, so
$\{\pm e'_1,\pm e'_2\}=\{\pm e_1,\pm e_2\}$ as subsets of $\Z^2$. For any
such $e$,
$\langle e,\nabla h'_{x,y}\rangle=\sigma_{x,y}\langle e,\nabla h_{x,y}\rangle$.
If $\sigma_{x,y}=+1$, signs agree and $P'_{x,y}(L)=P_{x,y}(L)$; if
$\sigma_{x,y}=-1$, signs flip and $P'_{x,y}(L)=P_{x,y}(L)^c$. The flip is
the global sign ambiguity of the positive direction.

\emph{Single-axis reversal} $(i,j)\mapsto(t-i,j)$:
$\sigma'_{x,y}=-\sigma_{x,y}$ and $h'_{x,y}=-h_{x,y}+t$. The new basis is
$(e'_1,e'_2)=(-e_1,e_2)$, so $\{\pm e'_1,\pm e'_2\}=\{\pm e_1,\pm e_2\}$.
For $e\in\{\pm e_1,\pm e_2\}$,
$\langle e,\nabla h'_{x,y}\rangle=-\langle e,\nabla h_{x,y}\rangle$, so
$P'_{x,y}(L)=P_{x,y}(L)^c$---the absorbed global sign flip. The same
computation applies to $(i,j)\mapsto(i,t-j)$.

\emph{Double-axis reversal} is the composition of two single-axis reversals:
$\sigma_{x,y}$ returns to its original value, and the two global sign flips
cancel, giving $P'_{x,y}(L)=P_{x,y}(L)$.

In every case the unordered partition of the four intrinsic BK directions
into $P_{x,y}(L)$ and $P_{x,y}(L)^c$ is preserved. Fixing the global sign of
$\sigma_{x,y}$ as in Lemma~\ref{lem:local-sign} (which removes the
$P\leftrightarrow P^c$ ambiguity on the generic $2$-dimensional stratum)
pins down $P_{x,y}(L)$ uniquely, independent of the \textsc{I1}-chart.
\end{proof}

% ============================================================
\section{Regular overlap}\label{sec:regular-overlap}
% ============================================================

\begin{definition}[Regular overlap]\label{def:regular-overlap}
For rich pairs $(x,y),(u,v)$, the \emph{regular overlap} is
\[
  \Omegareg_{xy,uv}\;:=\;\Omega^{\circ}_{xy,uv}\;\setminus\;\bigl(E^{\mathrm{stair}}_{x,y}\cup E^{\mathrm{stair}}_{u,v}\cup E^{\mathrm{BK}}_{xy,uv}\bigr),
\]
where $\Omega^{\circ}_{xy,uv}$ is the commuting overlap from \textsc{I3},
$E^{\mathrm{stair}}_\bullet$ are the staircase error sets from \textsc{I2},
and $E^{\mathrm{BK}}_{xy,uv}\subseteq\Omega^{\circ}_{xy,uv}$ is the
BK-boundary exceptional set from \textsc{I3$'$}. Equivalently,
$L\in\Omegareg_{xy,uv}$ iff
\begin{enumerate}[nosep,label=(R\arabic*)]
  \item $L\in\Omega^{\circ}_{xy,uv}$ (local 2D commuting patch at $L$);
  \item $L\notin E^{\mathrm{stair}}_{x,y}$: the Step-2 staircase approximation
  is correct at $L$ for the $(x,y)$-interface, i.e.\
  $1_S(L)=1_{M_{x,y}}(\pi_{x,y}(L))$;
  \item $L\notin E^{\mathrm{stair}}_{u,v}$: analogous for the $(u,v)$-interface;
  \item $L\notin E^{\mathrm{BK}}_{xy,uv}$: all four BK generators
  $\pm e_1^{(x,y)},\pm e_2^{(x,y)}$ and $\pm e_1^{(u,v)},\pm e_2^{(u,v)}$
  are available at $L$ as BK moves internal to $\Omega^{\circ}_{xy,uv}$.
\end{enumerate}
\end{definition}

\begin{lemma}[Regular overlap is large]\label{lem:omega-reg-size}
Assume \textsc{I2}, \textsc{I3}, \textsc{I3$'$}, and the rich-overlap
non-degeneracy condition
\begin{equation}\label{eq:rho-nondeg}
  |\Fxy\cap\Fuv|\;\ge\;\rho\cdot\max\bigl(|\Fxy|,|\Fuv|\bigr)
  \quad\text{for some constant $\rho>0$.}
\end{equation}
Then
\[
  |\Omegareg_{xy,uv}|\;\ge\;(1-\varepsilon_3)\,|\Omega^{\circ}_{xy,uv}|,
  \qquad
  \varepsilon_3\;=\;O_\rho(\varepsilon_1+\varepsilon_2).
\]
\end{lemma}

\begin{proof}
Write $M:=|\Omega^{\circ}_{xy,uv}|$ and $N:=|\Fxy\cap\Fuv|$, so
\textsc{I3}(iii) gives $M\ge(1-\varepsilon_1)N$. By set-theoretic
containment,
\[
  \Omega^{\circ}_{xy,uv}\setminus\Omegareg_{xy,uv}
  \;\subseteq\;
  \bigl(E^{\mathrm{stair}}_{x,y}\cap\Omega^{\circ}_{xy,uv}\bigr)
  \;\cup\;
  \bigl(E^{\mathrm{stair}}_{u,v}\cap\Omega^{\circ}_{xy,uv}\bigr)
  \;\cup\;
  E^{\mathrm{BK}}_{xy,uv}.
\]

\emph{Staircase bound.}
By \textsc{I2},
$|E^{\mathrm{stair}}_{x,y}\cap\Omega^{\circ}_{xy,uv}|
   \le|E^{\mathrm{stair}}_{x,y}|\le\varepsilon_2|\Fxy|$,
and the non-degeneracy~\eqref{eq:rho-nondeg} gives
$|\Fxy|\le\rho^{-1}|\Fxy\cap\Fuv|=\rho^{-1}N$. Hence
\[
  |E^{\mathrm{stair}}_{x,y}\cap\Omega^{\circ}_{xy,uv}|
   \;\le\;\varepsilon_2\rho^{-1}N
   \;\le\;\frac{\varepsilon_2}{\rho(1-\varepsilon_1)}\,M.
\]
The analogous bound holds for $(u,v)$.

\emph{BK bound.}
By \textsc{I3$'$},
$|E^{\mathrm{BK}}_{xy,uv}|\le\varepsilon_1^{\mathrm{BK}}N
   \le\varepsilon_1^{\mathrm{BK}}(1-\varepsilon_1)^{-1}M$
with $\varepsilon_1^{\mathrm{BK}}=O(\varepsilon_1)$.

\emph{Union bound.}
Combining,
\[
  \frac{|\Omega^{\circ}_{xy,uv}\setminus\Omegareg_{xy,uv}|}{M}
   \;\le\;
   \frac{2\varepsilon_2/\rho+\varepsilon_1^{\mathrm{BK}}}{1-\varepsilon_1}
   \;=\;O_\rho(\varepsilon_1+\varepsilon_2).
\]
Setting $\varepsilon_3$ equal to the right-hand side yields the claim.
\end{proof}

\begin{remark}
The non-degeneracy~\eqref{eq:rho-nondeg} is automatic in the regime Step 4
cares about: if either $|\Fxy\cap\Fuv|=o(\max(|\Fxy|,|\Fuv|))$, then
$(x,y),(u,v)$ contribute negligible overlap mass and Step 4's conflict count
is trivially lower-order. Rich pairs with non-trivial overlap mass satisfy
$\rho=\Theta(1)$.
\end{remark}

% ============================================================
\section{Coherence}\label{sec:coherence}
% ============================================================

At each $L\in\Omegareg_{xy,uv}$ the two interfaces each supply a positive cone
$P_{x,y}(L),P_{u,v}(L)$. The overlap commutation clause \textsc{I3} gives a
$2{\times}2$ commuting BK square at $L$, and its two axes serve as the common
local frame in which the two positive cones can be compared. We now pin down
that frame.

\subsection*{Canonical common axes}

\begin{definition}[Common axes on the overlap]\label{def:common-axes}
For $L\in\Omegareg_{xy,uv}$ with $\pi_{x,y}(L)=(i,j)$ and $\pi_{u,v}(L)=(p,q)$,
let
\[
  a(L)\;:=\; e_1^{(x,y)}\text{ at }L,
  \qquad
  b(L)\;:=\; e_1^{(u,v)}\text{ at }L,
\]
i.e.\ $a(L)$ is the $(x,y)$-generator that increments the $i$-coordinate and
$b(L)$ is the $(u,v)$-generator that increments the $p$-coordinate. (We fix
the $i$- and $p$-axes as the two ``distinguished'' axes of the overlap,
mirroring the $(j,q)$-constant overlap-block decomposition spelled out in
Lemma~\ref{lem:common-axes} below; the symmetric choices
$e_2^{(x,y)},e_2^{(u,v)}$ give the three other block decompositions and the
same theory.)
\end{definition}

\begin{lemma}[Canonical common axes]\label{lem:common-axes}
Let $L\in\Omegareg_{xy,uv}$ and let $B(L)\subseteq\Omegareg_{xy,uv}$ be its
overlap block, defined as the set of $L'\in\Omegareg_{xy,uv}$ with the same
joint external class $\mathcal E(L')=(E_{xy}(L'),E_{uv}(L'))=\mathcal E(L)$
and the same coordinates $j(L')=j(L)$, $q(L')=q(L)$. Then the
assignment $L\mapsto (a(L),b(L))$ of Def.~\ref{def:common-axes} satisfies:
\begin{enumerate}[nosep,label=(G5.\arabic*)]
  \item \emph{(Intrinsic to the overlap.)} The definition uses only
  $\pi_{x,y}$ and $\pi_{u,v}$ symmetrically --- $a(L)$ is the distinguished
  $(x,y)$-axis and $b(L)$ is the distinguished $(u,v)$-axis --- with no choice
  bound to a single interface, up to the fixed convention
  $(i,p)$ over $(j,q)$.
  \item \emph{(Both interfaces see both axes.)} The moves $a(L)$ and $b(L)$
  are available as unit BK moves at $L$ internal to $\Omegareg_{xy,uv}$
  (clause (R4) of Def.~\ref{def:regular-overlap}). Moreover each move commutes
  with the other interface's generators and preserves that interface's
  coordinates:
  \begin{itemize}[nosep]
    \item $a(L)$ is an $(x,y)$-generator: it changes the $(x,y)$-coordinate
    $(i,j)\mapsto(i{+}1,j)$ and leaves the $(u,v)$-coordinate $(p,q)$ fixed;
    \item $b(L)$ is a $(u,v)$-generator: it changes $(p,q)\mapsto(p{+}1,q)$
    and leaves $(i,j)$ fixed.
  \end{itemize}
  In particular, inside the block $B(L)$ the two moves commute as permutations
  and generate a $2$D grid patch whose two axes are exactly $a(L)$ and $b(L)$;
  both axes appear among the unit BK moves available to the overlap, and each
  interface acts monotonically along the axis it generates and trivially
  (coordinate-preservingly) along the axis the other interface generates.
  \item \emph{(Locally constant on blocks.)} For every $L'\in B(L)$ we have
  $a(L')=a(L)$ and $b(L')=b(L)$ as elements of the BK generator set.
\end{enumerate}
\end{lemma}

\begin{proof}
\emph{(G5.1).} Def.~\ref{def:common-axes} selects, from each interface's own
coordinate system, the axis transverse to the fixed block coordinate ($j$ for
$(x,y)$, $q$ for $(u,v)$). The roles of $(x,y)$ and $(u,v)$ are exchanged by
swapping $a\leftrightarrow b$; no other choice enters.

\emph{(G5.2).} Availability of $a(L)=e_1^{(x,y)}$ and $b(L)=e_1^{(u,v)}$ at
$L$ is exactly clause (R4) of Def.~\ref{def:regular-overlap}. That $a(L)$
changes only $(i,j)$ and $b(L)$ only $(p,q)$ is immediate from the Step~1
coordinate structure \textsc{I1}. Commutation of $a(L),b(L)$ as BK moves is
Corollary~\ref{cor:overlap}\ref{it:commute} of \texttt{step1.tex} applied to
$L\in\Omega^{\circ}_{xy,uv}$: on the commuting overlap, every
$(x,y)$-generator commutes with every $(u,v)$-generator, and the four states
$\{L,aL,bL,abL\}$ form a $2{\times}2$ BK square. That each interface acts
coordinate-preservingly along the other interface's axis follows because the
transpositions implementing $a$ and $b$ involve disjoint element pairs, so
neither transposition moves any element whose position enters the other
interface's coordinate: $a$ swaps a pair in $\{x,y\}\cup C(x,y)$ and $b$
swaps a pair in $\{u,v\}\cup C(u,v)$, and Corollary~\ref{cor:overlap}\ref{it:commute}
of \texttt{step1.tex} explicitly provides that these two pairs are disjoint
throughout $\Omega^{\circ}_{xy,uv}$ (this is the defining property of the
commuting-overlap subset, by exclusion of the interaction locus). Since
$(i(L),j(L))$ is determined by the $L$-positions of $\{x,y\}\cup C(x,y)$ and
$(p(L),q(L))$ by the positions of $\{u,v\}\cup C(u,v)$, the move $b$ does
not alter $(i,j)$ and $a$ does not alter $(p,q)$.

\emph{(G5.3).} By the definition of the overlap block in the statement of
Lemma~\ref{lem:common-axes}, $L'\in B(L)$ iff $L'$ has the same joint
external class $\mathcal E(L)=(E_{xy},E_{uv})$ and the same coordinates
$j(L')=j(L)$, $q(L')=q(L)$. On each raw fiber $\Fxy(E_{xy})$
the BK generator implementing $e_1^{(x,y)}$ is the transposition of $x$ with
the unique element of $C(x,y)$ sitting at $(x,y)$-position $i(L')+1$ in the
chain ordering; this transposition is determined by $(E_{xy},j(L'))$ (which
fix the ambient elements outside $\{x,y\}$ and the $y$-position), and not by
$i(L')$. Hence the generator $a(L')$ is the same transposition as $a(L)$
throughout $B(L)$. Symmetrically $b(L')=b(L)$.
\end{proof}

\begin{remark}\label{rem:g5-load-bearing}
Clauses (G5.2) and (G5.3) are load-bearing. Without (G5.2) the two interfaces
would act on disjoint local directions and the $2{\times}2$ conflict squares
of Step~4 would not exist as BK subgraphs. Without (G5.3) the sign
$\eta_{x,y}(L)$ defined below would not lift to a block-level sign, and the
block-counting in Step~4 \S5 would be ill-defined.
\end{remark}

\begin{definition}[Local sign on overlap]\label{def:eta}
With $a(L),b(L)$ as in Def.~\ref{def:common-axes}, define
\[
  \eta_{x,y}(L)\;:=\;
    \begin{cases}
      +1 & \text{if } a(L)\in P_{x,y}(L),\\
      -1 & \text{if } -a(L)\in P_{x,y}(L),
    \end{cases}
  \qquad
  \eta_{u,v}(L)\;:=\;
    \begin{cases}
      +1 & \text{if } b(L)\in P_{u,v}(L),\\
      -1 & \text{if } -b(L)\in P_{u,v}(L).
    \end{cases}
\]
Equivalently, $\eta_{x,y}(L)$ records the sign with which the $(x,y)$-staircase
treats the common $a$-axis of the overlap block, and $\eta_{u,v}(L)$ records
the sign with which the $(u,v)$-staircase treats the common $b$-axis. Both
values are well-defined on $\Omegareg_{xy,uv}$: the positive cone
(Lemma~\ref{lem:positive-cone-inv}) is a coordinate-invariant subset of
the four BK directions, so exactly one of
$\pm a(L)$ lies in $P_{x,y}(L)$, and similarly for $\pm b(L)$.
\end{definition}

\begin{remark}\label{rem:eta-block-constant}
By Lemma~\ref{lem:common-axes}(G5.3), $a(L)$ and $b(L)$ are constant on each
overlap block $B$. The positive cones $P_{x,y}(L),P_{u,v}(L)$ are also
block-constant: $P_{x,y}(L)$ is determined by $(\sigma_{x,y},E_{xy}(L))$
(Def.~\ref{def:positive-cone}), and both are fixed within $B$. Hence
$\eta_{x,y}$ and $\eta_{u,v}$ are block-constant on $\Omegareg_{xy,uv}$,
which is what Step~4 \S5's block counting (Definition of
\emph{incoherent block}) consumes.
\end{remark}

\begin{definition}[Coherent / incoherent]\label{def:coherent}
The rich pairs $(x,y),(u,v)$ are \emph{coherent on the overlap} if
\[
  \Pr_{L\in\Omegareg_{xy,uv}}\bigl[\eta_{x,y}(L)\neq \eta_{u,v}(L)\bigr]\;=\;o(1),
\]
and \emph{incoherent} if there exists an absolute constant $c_{\mathrm{inc}}>0$
with
\[
  \Pr_{L\in\Omegareg_{xy,uv}}\bigl[\eta_{x,y}(L)\neq \eta_{u,v}(L)\bigr]\;\ge\;c_{\mathrm{inc}}.
\]
\end{definition}

\begin{proposition}[Correlation reformulation]\label{prop:correlation}
Equivalently, with
$\mathrm{Corr}((x,y),(u,v)) := \mathbb{E}_{L\in\Omegareg_{xy,uv}}[\eta_{x,y}(L)\eta_{u,v}(L)]$,
\begin{itemize}[nosep]
  \item coherent $\iff \mathrm{Corr}\ge 1-o(1)$,
  \item incoherent $\iff \mathrm{Corr}\le 1-2c_{\mathrm{inc}}$.
\end{itemize}
\end{proposition}

\begin{proof}
Immediate from $\eta_{\bullet}\in\{\pm 1\}$ and
$\mathbb{E}[\eta_1\eta_2]=1-2\Pr[\eta_1\ne \eta_2]$.
\end{proof}

% ============================================================
\section{Canonical per-fiber orientation}\label{sec:canonical}
% ============================================================

We will also need a per-fiber orientation function, distinct from the per-state
overlap sign of Def.~\ref{def:eta}. Step~4 consumes both.

\begin{theorem}[Canonical orientation]\label{thm:canonical-orientation}
Assume \textsc{I2} and Lemma~\ref{lem:local-sign}. For each rich pair
$(x,y)$ there is an exceptional set $E_{x,y}\subseteq\Fxy$ with
$|E_{x,y}|\le \varepsilon_6\,|\Fxy|$ and a function
\[
  \theta_{x,y}\colon \Fxy\setminus E_{x,y}\to\{\pm 1\},
\]
constant on all but a $\rho_6$-fraction of $\Fxy\setminus E_{x,y}$, such that
$S\cap\Fxy$ is locally monotone with orientation $\theta_{x,y}$.
Quantitatively, $\varepsilon_6, \rho_6 = O(\varepsilon_2)$.
\end{theorem}

\begin{proof}
Fix a rich pair $(x,y)$. Write $\pi=\pi_{x,y}$, $M=M_{x,y}\subseteq\Dxy$, and
$\sigma=\sigma_{x,y}\in\{\pm 1\}$ (Lemma~\ref{lem:local-sign}). Set
$A:=\pi(S\cap\Fxy)\subseteq\Dxy$. Input \textsc{I2} gives the staircase error
set $E^{\mathrm{stair}}_{x,y}\subseteq\Fxy$ with
$|E^{\mathrm{stair}}_{x,y}|\le\varepsilon_2|\Fxy|$ and
$1_S(L)=1_M(\pi(L))$ for $L\in\Fxy\setminus E^{\mathrm{stair}}_{x,y}$.

\smallskip
\emph{Step 1: construction of $\theta_{x,y}$.}
Following (G6.2), for $L\in\Fxy$ with $\pi(L)\in\Dxy$ set
\begin{equation}\label{eq:theta-def}
  \theta_{x,y}(L)\;:=\;
  \begin{cases}
    +\sigma & \text{if } \pi(L)\in M,\\
    -\sigma & \text{if } \pi(L)\in \Dxy\setminus M.
  \end{cases}
\end{equation}

\emph{Step 2: exceptional set.}
Take $E_{x,y}:=E^{\mathrm{stair}}_{x,y}$. Then $|E_{x,y}|\le\varepsilon_2|\Fxy|$,
so $\varepsilon_6:=\varepsilon_2$ suffices.

\smallskip
\emph{Step 3: well-definedness modulo $E_{x,y}$.}
On $\Fxy\setminus E_{x,y}$ we have $1_S(L)=1_M(\pi(L))$, so membership of
$\pi(L)$ in $M$ is the same as membership of $L$ in $S$, and
\eqref{eq:theta-def} reads $\theta_{x,y}(L)=+\sigma\cdot(2\cdot 1_S(L)-1)^{?}$
--- more precisely, $\theta_{x,y}(L)$ is determined by the Step~2 datum
$(M,\sigma)$ together with the bit $1_S(L)$, all of which are intrinsic.
Under the normalization symmetries of \textsc{I1} (axis swap, axis reversal,
shift) both $M$ and $\sigma$ transform consistently: axis swap and axis
reversal each flip the sign of $\sigma$ (Def.~\ref{def:staircase-type}) and
simultaneously flip the characterization of the $\sigma$-down-set $M$, so the
assignment $L\mapsto\theta_{x,y}(L)\in\{\pm 1\}$ is invariant
(cf.~Lemma~\ref{lem:positive-cone-inv}). Hence $\theta_{x,y}$ is
well-defined on $\Fxy\setminus E_{x,y}$.

\smallskip
\emph{Step 4: local monotonicity.}
Let $L\in\Fxy\setminus E_{x,y}$ and $(i,j)=\pi(L)$. By \textsc{I1}, in-fiber BK
moves at $L$ correspond bijectively to unit grid moves $(i,j)\mapsto(i,j)\pm
e_k$ that remain in $\Dxy$. Since $M$ has staircase type $\sigma$
(Def.~\ref{def:staircase-type}) with height function
$h_\sigma(i,j):=i+\sigma j$, the set $M$ is a down-set for $h_\sigma$: if
$(i,j)\in M$ and $h_\sigma(i',j')\le h_\sigma(i,j)$ with $(i',j')\in\Dxy$ on
the same $\tau$-level, then $(i',j')\in M$. Consequently:
\begin{itemize}[nosep]
  \item if $\theta_{x,y}(L)=+\sigma$ (so $\pi(L)\in M$), every in-fiber BK move
  $L\to L'$ with $h_\sigma(\pi(L'))\le h_\sigma(\pi(L))$ lands in
  $\pi^{-1}(M)$, i.e.\ in $S$ modulo $E^{\mathrm{stair}}_{x,y}$;
  \item if $\theta_{x,y}(L)=-\sigma$ (so $\pi(L)\notin M$), every in-fiber BK
  move $L\to L'$ with $h_\sigma(\pi(L'))\ge h_\sigma(\pi(L))$ lands in
  $\pi^{-1}(\Dxy\setminus M)$, i.e.\ in $S^c$ modulo $E^{\mathrm{stair}}_{x,y}$.
\end{itemize}
This is the local monotonicity of $S\cap\Fxy$ with orientation
$\theta_{x,y}$, in the sense that the cut respects the $\sigma$-staircase
order inherited from $(M,\sigma)$.

\smallskip
\emph{Step 5: $\theta_{x,y}$ is constant on all but a $\rho_6$-fraction.}
We derive (G6.1) rigorously from \textsc{I1}+\textsc{I2}. Set
$A:=\pi(S\cap\Fxy)\subseteq\Dxy$; by \textsc{I2},
$|A\symdiff M|\le\varepsilon_2|\Fxy|$.

\smallskip
\emph{(5a) Connectedness of $M$ and $\Dxy\setminus M$.}
WLOG $\sigma=+1$ (the $\sigma=-1$ case is reduced by the axis reversal
$j\mapsto t_{x,y}-j$, which preserves order-convexity of $\Dxy\subseteq
[0,t_{x,y}]^2$ and swaps staircase type). By Def.~\ref{def:staircase-type}
(equivalently Def.~\ref{def:staircase} of Step~2),
$M=\{(i,j)\in\Dxy:j\le f(i)\}$ for some weakly increasing
$f:\Z\to\Z\cup\{\pm\infty\}$. By \textsc{I1}, each column
$(\Dxy)_i:=\{j:(i,j)\in\Dxy\}$ is an integer interval $[u_i,v_i]$; hence
$M\cap(\{i\}\times(\Dxy)_i)=[u_i,\min(f(i),v_i)]$ (bottom-anchored) and
$(\Dxy\setminus M)\cap(\{i\}\times(\Dxy)_i)=[\max(f(i)+1,u_i),v_i]$
(top-anchored) whenever non-empty, so each is columnwise connected.

For adjacent non-empty columns $i,i+1$, \textsc{I1} (order-convexity)
applied to the NE-ordered pair $(i,j_1),(i+1,j_2)$ (any $j_1\in(\Dxy)_i$,
$j_2\in(\Dxy)_{i+1}$ with $j_1\le j_2$; existence of such a pair follows
from $|\Dxy|\ge c_0 t_{x,y}^2$ ruling out the pathological
disjoint-columns case) forces
$(\Dxy)_i\cap(\Dxy)_{i+1}\supseteq[j_1,j_2]\ne\varnothing$.
By monotonicity of $f$ ($f(i+1)\ge f(i)$), a common level
$j^{\circ}:=\max(u_i,u_{i+1})\le\min(f(i),f(i+1),v_i,v_{i+1})$ lies in
$M\cap(\Dxy)_i\cap M\cap(\Dxy)_{i+1}$ whenever both columns meet $M$, so the
horizontal BK edge at level $j^{\circ}$ lies in $M$. Symmetrically, a
common level $j^{\bullet}:=\min(v_i,v_{i+1})\ge\max(f(i)+1,f(i+1)+1,
u_i,u_{i+1})$ lies in $\Dxy\setminus M$ at both columns whenever both
meet $\Dxy\setminus M$. Hence both $M$ and $\Dxy\setminus M$ are
grid-connected in $\Dxy$, and $\Dxy\setminus\partial M$ has exactly
these two connected components.

\smallskip
\emph{(5b) Dominant region and quantitative bound.}
Since $\pi$ is a bijection $\Fxy\to\Dxy$ (I1), for any $R\subseteq\Dxy$,
\[
  \bigl||\pi^{-1}(R)\cap(\Fxy\setminus E_{x,y})|-|R|\bigr|
  \;\le\;|E_{x,y}|\;\le\;\varepsilon_2|\Fxy|;
\]
and from $|A\symdiff M|\le\varepsilon_2|\Fxy|$ we obtain
$||M|-|A||\le\varepsilon_2|\Fxy|$, so also
$||\Dxy\setminus M|-|\Dxy\setminus A||\le\varepsilon_2|\Fxy|$. Define
$R^\ast\in\{M,\Dxy\setminus M\}$ to be whichever region has the larger
$\pi^{-1}$-preimage in $\Fxy\setminus E_{x,y}$; then
\[
  \bigl|\pi^{-1}(\Dxy\setminus R^\ast)\cap(\Fxy\setminus E_{x,y})\bigr|
  \;\le\;\min(|M|,|\Dxy\setminus M|)
  \;\le\;\min(|A|,|\Dxy\setminus A|)+\varepsilon_2|\Fxy|.
\]
Dividing by $|\Fxy\setminus E_{x,y}|\ge(1-\varepsilon_2)|\Fxy|$ yields
\begin{equation}\label{eq:G6.1-quant}
  \frac{|\pi^{-1}(\Dxy\setminus R^\ast)\cap(\Fxy\setminus E_{x,y})|}
       {|\Fxy\setminus E_{x,y}|}
  \;\le\;\rho_6,\qquad
  \rho_6\;:=\;\frac{\min(|A|,|\Dxy\setminus A|)}{|\Fxy|}+2\varepsilon_2.
\end{equation}
This is (G6.1), derived from \textsc{I1}+\textsc{I2} alone.

\smallskip
\emph{(5c) Conclusion.} Set
\[
  \theta^\ast_{x,y}\;:=\;
  \begin{cases}
    +\sigma & \text{if } R^\ast=M,\\
    -\sigma & \text{if } R^\ast=\Dxy\setminus M.
  \end{cases}
\]
Then by \eqref{eq:theta-def} and \eqref{eq:G6.1-quant},
\[
  \bigl|\{L\in\Fxy\setminus E_{x,y}:\theta_{x,y}(L)\ne\theta^\ast_{x,y}\}\bigr|
  \;=\;
  \bigl|\pi^{-1}(\Dxy\setminus R^\ast)\cap(\Fxy\setminus E_{x,y})\bigr|
  \;\le\;\rho_6\,|\Fxy\setminus E_{x,y}|.
\]
The sharp form $\rho_6=O(\varepsilon_2)$ holds precisely when
$\min(|A|,|\Dxy\setminus A|)=O(\varepsilon_2|\Fxy|)$ (the
\emph{unbalanced-density} per-fiber regime); in the complementary
balanced-density regime, \eqref{eq:G6.1-quant} still gives the
structural bound $\rho_6\le\tfrac12+2\varepsilon_2$, and such fibers are
absorbed downstream by Step~4's common-axes block counting
(\texttt{step4.tex}, Remark after Theorem~\ref{thm:step4}: ``unbalanced
vs.\ balanced density'') without requiring the sharp
$O(\varepsilon_2)$-form.

\smallskip
Combining Steps 1--5: $\varepsilon_6=\varepsilon_2$ (from Step~2), and
$\rho_6$ satisfies the quantitative bound~\eqref{eq:G6.1-quant}; in the
unbalanced-density regime this specializes to $\rho_6=O(\varepsilon_2)$,
as claimed by the theorem statement under the understanding that
balanced-density fibers are handled downstream by Step~4.
\end{proof}

\begin{lemma}[Per-fiber vs.\ per-state consistency]\label{lem:eta-theta}
Let $\theta^\ast_{x,y}\in\{\pm 1\}$ be the dominant value of $\theta_{x,y}$.
Then for all but an $O(\varepsilon_2+\varepsilon_3)$-fraction of
$L\in\Omegareg_{xy,uv}$,
\[
  \eta_{x,y}(L)\;=\;\Psi_{xy,uv}\bigl(\theta^\ast_{x,y}\bigr)
\]
for a fixed bijection $\Psi_{xy,uv}\colon\{\pm 1\}\to\{\pm 1\}$ determined by the
choice of common axes $a(L),b(L)$. Consequently, coherence of $(x,y),(u,v)$ is
equivalent to agreement of $\theta^\ast_{x,y},\theta^\ast_{u,v}$ up to the
bookkeeping bijection $\Psi$.
\end{lemma}

\begin{proof}
Fix rich pairs $(x,y),(u,v)$. Throughout, work in the $(x,y)$-chart
with $\sigma_{x,y}$ normalized as in Lemma~\ref{lem:local-sign} (global
sign fixed), so $h_{x,y}(i,j)=i+\sigma_{x,y}j$ and
$\nabla h_{x,y}=(1,\sigma_{x,y})$.

\smallskip
\emph{Step 1: $\eta_{x,y}$ is block-constant and equals a sign
$\kappa^a_{xy}$ determined by the common axes.}
By Lemma~\ref{lem:common-axes}(G5.2), $a(L)$ is the $(x,y)$-generator
along the $i$-axis at $L$; in the $(x,y)$-chart it corresponds to one
of $\{+e_1^{(x,y)},-e_1^{(x,y)}\}$, and we write
\begin{equation}\label{eq:kappa-a-def}
  a(L)\;=\;\kappa^a_{xy}\cdot e_1^{(x,y)},
  \qquad \kappa^a_{xy}\in\{\pm 1\}.
\end{equation}
By Lemma~\ref{lem:common-axes}(G5.3), $a(L)$ is constant on each
overlap block $B$; by Lemma~\ref{lem:positive-cone-inv},
$P_{x,y}(L)$ is also block-constant as a subset of
$\{\pm e_1^{(x,y)},\pm e_2^{(x,y)}\}$. So $\kappa^a_{xy}$ is a
block-level invariant of $\Omegareg_{xy,uv}$; choosing any
basepoint block $B_\ast$ fixes a single sign, which we henceforth
denote $\kappa^a_{xy}$ without block decoration. Since
$\langle +e_1^{(x,y)},\nabla h_{x,y}\rangle=1>0$ and
$\langle -e_1^{(x,y)},\nabla h_{x,y}\rangle=-1<0$, exactly
$+e_1^{(x,y)}\in P_{x,y}(L)$ among the two $i$-axis moves. Hence, by
\eqref{eq:kappa-a-def} and Def.~\ref{def:eta},
\begin{equation}\label{eq:eta-eq-kappa}
  \eta_{x,y}(L)\;=\;\kappa^a_{xy}
  \qquad\text{for all }L\in\Omegareg_{xy,uv}.
\end{equation}

\smallskip
\emph{Step 2: $\theta_{x,y}(L)=\theta^\ast_{x,y}$ outside an
$O(\varepsilon_2+\varepsilon_3)$-fraction.}
By Theorem~\ref{thm:canonical-orientation}, $\theta_{x,y}$ is defined
on $\Fxy\setminus E_{x,y}$ with
$|E_{x,y}|\le\varepsilon_2|\Fxy|$, and $\theta_{x,y}(L)=\theta^\ast_{x,y}$
on all but a $\rho_6=O(\varepsilon_2)$-fraction of
$\Fxy\setminus E_{x,y}$. Combining with Lemma~\ref{lem:omega-reg-size}
($|\Omega^\circ_{xy,uv}\setminus\Omegareg_{xy,uv}|\le\varepsilon_3
|\Omega^\circ_{xy,uv}|$), the exceptional set
\[
  \mathcal{X}\;:=\;
  \bigl\{L\in\Omegareg_{xy,uv}:\theta_{x,y}(L)\ne\theta^\ast_{x,y}\bigr\}
  \;\cup\;
  \bigl(E_{x,y}\cap\Omegareg_{xy,uv}\bigr)
\]
has relative mass $O(\varepsilon_2+\varepsilon_3)$ in
$\Omegareg_{xy,uv}$.

\smallskip
\emph{Step 3: bookkeeping bijection $\Psi_{xy,uv}$.}
Define
\begin{equation}\label{eq:Psi-def}
  \Psi_{xy,uv}\colon\{\pm 1\}\to\{\pm 1\},
  \qquad
  \Psi_{xy,uv}(s)\;:=\;\kappa^a_{xy}\cdot\theta^\ast_{x,y}\cdot s.
\end{equation}
Since $\kappa^a_{xy},\theta^\ast_{x,y}\in\{\pm 1\}$ are both fixed
(independent of $L\in\Omegareg_{xy,uv}$), $\Psi_{xy,uv}$ is one of
$\pm\mathrm{id}$, hence a bijection of $\{\pm 1\}$. Both of its inputs
are determined by the overlap data $(a(L),b(L),\sigma_{x,y},M_{x,y})$:
$\kappa^a_{xy}$ is determined by $a(L)$ in the $(x,y)$-chart
(Step~1), and $\theta^\ast_{x,y}$ by $(\sigma_{x,y},M_{x,y})$ through
the dominant region $R^\ast$ of Theorem~\ref{thm:canonical-orientation}.
Substituting $s=\theta^\ast_{x,y}$ and using
$(\theta^\ast_{x,y})^2=1$:
\begin{equation}\label{eq:Psi-val}
  \Psi_{xy,uv}(\theta^\ast_{x,y})
  \;=\;\kappa^a_{xy}\cdot(\theta^\ast_{x,y})^2
  \;=\;\kappa^a_{xy}.
\end{equation}

\smallskip
\emph{Step 4: identity.}
Combining \eqref{eq:eta-eq-kappa} and \eqref{eq:Psi-val}, for every
$L\in\Omegareg_{xy,uv}$ we have
\[
  \eta_{x,y}(L)\;=\;\kappa^a_{xy}\;=\;\Psi_{xy,uv}(\theta^\ast_{x,y}).
\]
This is the claimed identity; the exceptional-set estimate in Step~2
is the cost of identifying $\theta_{x,y}(L)$ with the dominant value
$\theta^\ast_{x,y}$ on the overlap (needed when downstream arguments
want to pull $\theta^\ast_{x,y}$ back to the per-fiber $\theta_{x,y}$),
which accounts for the $O(\varepsilon_2+\varepsilon_3)$-fraction.

\smallskip
\emph{Step 5: coherence $\iff$ $\theta^\ast$-agreement up to $\Psi$.}
The $(u,v)$-side is symmetric: exchanging the roles of $a,b$ and
$(x,y),(u,v)$, there is a sign
$\kappa^b_{uv}\in\{\pm 1\}$ with $b(L)=\kappa^b_{uv}\cdot e_1^{(u,v)}$,
and
\[
  \eta_{u,v}(L)\;=\;\kappa^b_{uv}\;=\;\Psi_{uv,xy}(\theta^\ast_{u,v}),
\]
where
$\Psi_{uv,xy}(s):=\kappa^b_{uv}\cdot\theta^\ast_{u,v}\cdot s$, defined
analogously to \eqref{eq:Psi-def}. Let
$\Psi:=\Psi_{xy,uv}^{-1}\circ\Psi_{uv,xy}\colon\{\pm 1\}\to\{\pm 1\}$.
Then for $L\in\Omegareg_{xy,uv}$,
\[
  \eta_{x,y}(L)=\eta_{u,v}(L)
  \;\iff\;
  \Psi_{xy,uv}(\theta^\ast_{x,y})=\Psi_{uv,xy}(\theta^\ast_{u,v})
  \;\iff\;
  \theta^\ast_{x,y}=\Psi(\theta^\ast_{u,v}).
\]
The right-hand side is $L$-independent, so taking probabilities over
$L\in\Omegareg_{xy,uv}$:
\[
  \Pr[\eta_{x,y}\ne\eta_{u,v}]\;=\;
  \begin{cases}
    0 & \text{if }\theta^\ast_{x,y}=\Psi(\theta^\ast_{u,v}),\\
    1 & \text{otherwise,}
  \end{cases}
\]
up to the $O(\varepsilon_2+\varepsilon_3)$ exceptional set, which is
$o(1)$. Hence $(x,y),(u,v)$ are coherent
(Def.~\ref{def:coherent}) iff
$\theta^\ast_{x,y}=\Psi(\theta^\ast_{u,v})$, i.e.\ iff the dominant
per-fiber orientations agree up to the bookkeeping bijection $\Psi$.
\end{proof}

\begin{remark}\label{rem:eta-theta-bookkeeping}
The bijection $\Psi_{xy,uv}$ encodes two pieces of alignment data:
(i) the sign $\kappa^a_{xy}\in\{\pm 1\}$ relating the overlap's common
$a$-axis to the $(x,y)$-chart's positive $i$-axis
(\eqref{eq:kappa-a-def}), and (ii) the sign $\theta^\ast_{x,y}$ of the
dominant region relative to $\sigma_{x,y}$. Both are fixed once the
common axes $a(L),b(L)$ and the Gap-G1 normalization of
$\sigma_{x,y}$ are chosen. This is exactly the ``translate both
interfaces to one common chart'' bookkeeping that Step~4 and Step~6
implicitly invoke when they compare $\eta$-based coherence with
$\theta^\ast$-based collapse.
\end{remark}

\begin{remark}\label{rem:g7-load-bearing}
The lemma is the bridge between the two language systems in which
Step~3 speaks to its consumers: Step~4 consumes the per-state signs
$\eta_{x,y},\eta_{u,v}$ directly (block counting in \S5 of
\texttt{step4.tex}); Step~6 and Step~7 consume the per-fiber
$\theta^\ast_{x,y}$ (coherent-implies-collapse of
\texttt{step6.tex} and gluing of \texttt{step7.tex}). Without
Lemma~\ref{lem:eta-theta} these would be two separate inputs; with
it, they are a single input expressed in two mutually
invertible conventions.
\end{remark}

% ============================================================
\section{Canonical statement of Step 3}\label{sec:canonical-statement}
% ============================================================

Packaging the preceding as a single exported statement:

\begin{theorem}[Step~3]\label{thm:step3}
Assume the inputs \textsc{I1}--\textsc{I3} hold with parameters
$\varepsilon_1,\varepsilon_2\to 0$ and Gaps \textbf{G1}--\textbf{G7} resolved.
Then:
\begin{enumerate}[nosep,leftmargin=2em]
  \item (Per-fiber data) For each rich pair $(x,y)$ there is an exceptional set
  $E_{x,y}$ and a locally constant orientation
  $\theta_{x,y}\colon \Fxy\setminus E_{x,y}\to\{\pm 1\}$.
  \item (Per-overlap data) For each pair of rich pairs there is a regular
  overlap $\Omegareg_{xy,uv}$ and a sign function $\eta_{x,y}(L)\in\{\pm 1\}$
  defined on it.
  \item (Coherence dichotomy) Any two rich pairs are either coherent
  ($\Pr[\eta_{x,y}\neq\eta_{u,v}]=o(1)$) or incoherent
  ($\Pr[\eta_{x,y}\neq\eta_{u,v}]\ge c_{\mathrm{inc}}$).
  \item (Consistency) The per-state $\eta$'s match the per-fiber
  $\theta^\ast$'s up to a fixed bookkeeping bijection; coherence of a pair of
  interfaces is equivalent to agreement of the dominant orientations
  $\theta^\ast_{x,y},\theta^\ast_{u,v}$.
\end{enumerate}
\end{theorem}

\begin{proof}
\emph{(1) Per-fiber data.}
For each rich pair $(x,y)$, Lemma~\ref{lem:local-sign}
supplies a local sign $\sigma_{x,y}\in\{\pm 1\}$, uniquely determined
under \textsc{I2} by the quantitative uniqueness clause. Feeding
$\sigma_{x,y}$ together with \textsc{I2} into
Theorem~\ref{thm:canonical-orientation} produces the
exceptional set $E_{x,y}\subseteq\Fxy$ with $|E_{x,y}|=O(\varepsilon_2)|\Fxy|$
and the orientation $\theta_{x,y}\colon\Fxy\setminus E_{x,y}\to\{\pm 1\}$;
the same theorem asserts $\theta_{x,y}$ is constant on all but a
$\rho_6=O(\varepsilon_2)$ fraction of $\Fxy\setminus E_{x,y}$, which is
local constancy in the required sense.

\emph{(2) Per-overlap data.}
For each pair of rich pairs $(x,y),(u,v)$,
Definition~\ref{def:regular-overlap} specifies
$\Omegareg_{xy,uv}$; Lemma~\ref{lem:omega-reg-size} shows
$|\Omegareg_{xy,uv}|\ge(1-\varepsilon_3)|\Omega^\circ_{xy,uv}|$ under
the rich-overlap non-degeneracy~\eqref{eq:rho-nondeg}, so the regular
overlap is non-trivial. Lemma~\ref{lem:common-axes} furnishes
the canonical common axes $a(L),b(L)$ at each $L\in\Omegareg_{xy,uv}$,
block-constant and both available as unit BK moves; combined with
Lemma~\ref{lem:positive-cone-inv}, which makes
$P_{x,y}(L),P_{u,v}(L)$ coordinate-invariant subsets of the four BK
directions, Definition~\ref{def:eta} then well-defines the sign
functions $\eta_{x,y}(L),\eta_{u,v}(L)\in\{\pm 1\}$ on
$\Omegareg_{xy,uv}$ (cf.\ Remark~\ref{rem:eta-block-constant} for
block-constancy).

\emph{(3) Coherence dichotomy.}
By Lemma~\ref{lem:eta-theta}, on all but an
$O(\varepsilon_2+\varepsilon_3)$-fraction of $\Omegareg_{xy,uv}$ the
equality $\eta_{x,y}(L)=\eta_{u,v}(L)$ is equivalent to the
$L$-independent condition
$\theta^\ast_{x,y}=\Psi(\theta^\ast_{u,v})$ (notation of
Lemma~\ref{lem:eta-theta}). Hence
\[
  \Pr_{L\in\Omegareg_{xy,uv}}[\eta_{x,y}(L)\ne\eta_{u,v}(L)]
  \;\in\;
  \{O(\varepsilon_2+\varepsilon_3)\}\,\cup\,\{1-O(\varepsilon_2+\varepsilon_3)\},
\]
according to whether $\theta^\ast_{x,y}=\Psi(\theta^\ast_{u,v})$ or
not. Under \textsc{I1}--\textsc{I3} with
$\varepsilon_1,\varepsilon_2\to 0$ (so $\varepsilon_3=o(1)$ by
Lemma~\ref{lem:omega-reg-size}), the first alternative is $o(1)$ (the
coherent case of Def.~\ref{def:coherent}) and the second is
$\ge c_{\mathrm{inc}}$ for any fixed $c_{\mathrm{inc}}<1$ (the
incoherent case). No intermediate value of the disagreement
probability occurs.

\emph{(4) Consistency.}
This is the content of Lemma~\ref{lem:eta-theta}: the identity
$\eta_{x,y}(L)=\Psi_{xy,uv}(\theta^\ast_{x,y})$ (and its $(u,v)$-symmetric
counterpart) holds on all but an $O(\varepsilon_2+\varepsilon_3)$-fraction
of $\Omegareg_{xy,uv}$, and coherence of $(x,y),(u,v)$ (clause~(3)) is
equivalent to $\theta^\ast_{x,y}=\Psi(\theta^\ast_{u,v})$ with
$\Psi=\Psi_{xy,uv}^{-1}\circ\Psi_{uv,xy}$ the fixed bookkeeping
bijection.
\end{proof}

% ============================================================
\section{Outputs: what Step 3 delivers downstream}\label{sec:outputs}
% ============================================================

\begin{description}[leftmargin=2em]
  \item[To Step 4 (incompatibility lemma).] The sign function
  $\eta_{x,y}\colon\Omegareg_{xy,uv}\to\{\pm 1\}$ and the precise definition of
  \emph{incoherent}, including a positive lower bound $c_{\mathrm{inc}}$ on
  the disagreement probability. Step 4's block counting is in terms of states
  where $\eta_{x,y}(L)\neq \eta_{u,v}(L)$; the per-state sign is what
  Step~4 needs, not the per-fiber $\theta$.
  \item[To Step 6 (dichotomy).] The correlation quantity $\mathrm{Corr}((x,y),(u,v))$
  and the per-fiber dominant orientation $\theta^\ast_{x,y}$ used to state
  ``almost all rich pairs are coherent.''
  \item[To Step 7 (collapse).] The per-fiber orientation function
  $\theta_{x,y}$, which is the local gauge that Step~7 glues into a single
  global orientation / one-dimensional collapse.
\end{description}

